\section{有噪信道编码}

\begin{remarkBox}{本章重点}
第六章讨论信道容量，回答“信道最多能传多少信息”；本章进一步讨论有噪信道中怎样通过增加冗余提高可靠性。学习主线是：先理解信道编码和译码的基本模型，再掌握使平均错误率最小的 MAP 准则和等概输入下常用的 ML 准则，随后用汉明距离分析线性分组码的纠错能力，最后理解有噪信道编码定理给出的可靠传输极限。
\end{remarkBox}

有噪信道编码的核心思想不是消除信道噪声，而是在发送前主动加入有规律的冗余，使接收端即使收到被扰乱的序列，也能根据这种规律尽可能恢复原消息。容量给出了可靠通信的上限，编码和译码则回答如何靠近这个上限。

\subsection{有噪信道编码概述}

信道编码的目的是提高传输可靠性。发送端把信源输出的消息序列映射为满足一定数学规律的码字，码字中除了原始信息，还包含冗余符号；接收端再利用同样的规律进行译码，恢复原来的消息序列。

\begin{definitionBox}{信道编码与信道译码}
\textbf{信道编码}是按一定规则给信源输出序列增加冗余符号，使其变成满足一定数学规律的码序列或码字，再经信道传输。

\textbf{信道译码}是按与编码器对应的数学规律处理接收序列，去掉冗余影响并恢复信源消息序列。
\end{definitionBox}

设消息集合为 $\{1,2,\dots,M\}$，编码器把第 $i$ 个消息映射为码字 $\mathbf{c}_i$。若码长为 $n$，则信道输入为 $X^n$，信道输出或接收序列为 $Y^n$，译码器根据 $Y^n$ 输出恢复消息 $V$。这里的 $M$ 表示可区分的消息数，$n$ 表示每个码字占用的信道符号数。

\begin{formulaBox}{信息传输速率}
一个 $(M,n)$ 码的信息传输速率为
\[
R=\frac{H(X)}{n}.
\]
当 $M$ 个消息等概率时，
\[
R=\frac{\log M}{n}.
\]
若为二进制 $(n,k)$ 线性分组码，则 $M=2^k$，所以
\[
R=\frac{\log_2 2^k}{n}=\frac{k}{n},
\]
此时 $R$ 也称为码率或编码效率。
\end{formulaBox}

译码方式通常分为硬判决和软判决。硬判决先对信道传输符号作判决，再进行信道译码；软判决则把信道传输符号判决和信道译码合在一起完成。一般来说，软判决保留了更多接收信息，但实现也更复杂。

\subsection{判决与译码规则}

先看单符号译码。设信道输入字母表为
\[
A=\{a_1,a_2,\dots,a_r\},
\]
输出字母表为
\[
B=\{b_1,b_2,\dots,b_s\}.
\]
译码规则就是给每个输出符号指定一个输入符号作为判决结果。

\begin{definitionBox}{单符号判决规则}
译码或判决规则是一个单值函数
\[
F(y=b_j)=a_i,
\]
表示接收到 $b_j$ 时判定发送符号为 $a_i$。每一个信道输出都必须有且只有一个对应的判决输入。
\end{definitionBox}

如果接收到 $b_j$ 后按照规则判为 $a^*$，则只有实际发送 $a^*$ 时译码正确。对应的条件错误率为
\[
\sum_{a_i\ne a^*}P_{X\mid Y}(x=a_i\mid y=b_j).
\]
对所有输出加权平均，得到平均错误率
\[
P_E=\sum_j P_Y(y=b_j)\sum_{a_i\ne a^*}P_{X\mid Y}(x=a_i\mid y=b_j)
=\sum_y p(y)\sum_{x\ne x^*}p(x\mid y).
\]
由于每个 $y$ 只保留一个判决 $x^*$，平均正确率为
\[
1-P_E=\sum_y p(x^*,y),
\]
所以
\[
P_E=1-\sum_y p(x^*,y).
\]
做译码题时，通常就是为每一列输出选择一个输入，使被选中的联合概率之和尽可能大。

\begin{example}
\textbf{判决规则与平均错误率}：设
\[
[P_X]=[0.4\quad 0.6],
\qquad
[P_{Y\mid X}]=
\begin{bmatrix}
0.8&0.2\\
0.1&0.9
\end{bmatrix}.
\]
求四种判决规则的平均错误率，并判断最优规则。

\noindent \textbf{解题流程}：\\[2pt]
步骤一：由先验概率和转移概率得到联合概率矩阵。\\
步骤二：列出每个输出对应输入的判决规则。\\
步骤三：把规则选中的联合概率相加得到正确率，再用 $1-$正确率得到错误率。

\textbf{解}：联合概率矩阵为
\[
[P_{XY}]=
\begin{bmatrix}
0.32&0.08\\
0.06&0.54
\end{bmatrix}.
\]
四种规则分别为
\[
F_1(b_1)=a_1,\quad F_1(b_2)=a_1,
\]
\[
F_2(b_1)=a_2,\quad F_2(b_2)=a_2,
\]
\[
F_3(b_1)=a_1,\quad F_3(b_2)=a_2,
\]
\[
F_4(b_1)=a_2,\quad F_4(b_2)=a_1.
\]
对应平均错误率为
\[
P_e(F_1)=0.6,
\quad
P_e(F_2)=0.4,
\quad
P_e(F_3)=0.14,
\quad
P_e(F_4)=0.86.
\]
因此 $F_3$ 最好，$F_4$ 最差。
\end{example}

\subsection{最佳判决与译码准则}

判决规则很多，但课程中最常用的是最大后验概率准则和最大似然准则。前者利用先验概率和信道统计特性，直接最小化平均错误率；后者只比较转移概率，在输入等概或先验未知时特别常用。

\begin{definitionBox}{最大后验概率准则（MAP）}
对所有 $i$，若
\[
p(x=a^*\mid y)\ge p(x=a_i\mid y),
\]
则选择
\[
F(y)=a^*.
\]
这称为最大后验概率准则。MAP 准则把后验概率最大的输入符号作为译码输出，是使平均错误率最小的判决准则。
\end{definitionBox}

由贝叶斯公式，MAP 判决也可写成似然比检验：
\[
\frac{p(y\mid x=a^*)}{p(y\mid x=a_i)}
\ge
\frac{p(x=a_i)}{p(x=a^*)}.
\]
因此 MAP 既看信道转移概率，也看输入符号先验概率。等价地，对固定输出 $b_j$，MAP 可以逐列比较联合概率
\[
P(a^*,b_j)\ge P(a_i,b_j).
\]

\begin{definitionBox}{最大似然准则（ML）}
若对所有 $i$ 都有
\[
p(y\mid x=a^*)\ge p(y\mid x=a_i),
\]
则选择
\[
F(y)=a^*.
\]
这称为最大似然准则。ML 准则按最大转移概率作判决；当输入符号等概时，ML 准则与 MAP 准则等价。
\end{definitionBox}

比较两种准则时，可以记住一个做题入口：MAP 先把转移概率矩阵每一行乘以对应先验概率，得到联合概率矩阵，然后每列选最大；ML 直接在转移概率矩阵中每列选最大。若输入等概，联合概率矩阵只是转移概率矩阵整体乘以同一个常数，因此两种准则给出相同判决。

\textbf{MAP/ML 做题流程}：自学时可以按下面顺序判断，不必一开始就套复杂公式。
\begin{enumerate}
  \item 先看题目是否给出输入先验概率 $p(a_i)$。若给出且不等概，优先使用 MAP；若输入等概或先验未知，通常使用 ML。
  \item MAP 题先算联合概率
  \[
  p(a_i,b_j)=p(a_i)p(b_j\mid a_i),
  \]
  然后对每个输出 $b_j$ 在同一列中选最大的联合概率。
  \item ML 题直接对每个输出 $b_j$ 比较 $p(b_j\mid a_i)$，也就是在转移概率矩阵同一列中选最大元素。
  \item 平均正确率等于“每列被选中项”的联合概率之和，平均错误率等于 $1-$平均正确率。
  \item 若处理的是码字序列，在无记忆信道中把似然分解为 $\prod_i p(y_i\mid x_i)$；在 $p\le1/2$ 的 BSC 中，可进一步转化为最小汉明距离译码。
\end{enumerate}

\begin{example}
\textbf{例7.2：MAP 与 ML 判决}：设信道输入 $X$ 取值为 $(a_1,a_2,a_3)$，概率分别为 $1/2,1/4,1/4$。信道输出 $Y$ 取值为 $(b_1,b_2,b_3)$，转移概率矩阵为
\[
P=
\begin{bmatrix}
0.5&0.3&0.2\\
0.2&0.3&0.5\\
0.2&0.4&0.4
\end{bmatrix}.
\]
求 MAP 和 ML 准则下的判决函数及错误率。

\noindent \textbf{解题流程}：\\[2pt]
步骤一：ML 直接对转移概率矩阵逐列找最大。\\
步骤二：MAP 先求联合概率矩阵，再逐列找最大。\\
步骤三：被选中联合概率之和为平均正确率，错误率为 $1-$正确率。

\textbf{解}：ML 准则逐列比较 $P(y\mid x)$，得到
\[
F_{\mathrm{ML}}(b_1)=a_1,
\quad
F_{\mathrm{ML}}(b_2)=a_3,
\quad
F_{\mathrm{ML}}(b_3)=a_2.
\]
联合概率矩阵为
\[
[P(a_i,b_j)]=
\begin{bmatrix}
1/4&3/20&1/10\\
1/20&3/40&1/8\\
1/20&1/10&1/10
\end{bmatrix}.
\]
MAP 准则逐列比较联合概率，得到
\[
F_{\mathrm{MAP}}(b_1)=a_1,
\quad
F_{\mathrm{MAP}}(b_2)=a_1,
\quad
F_{\mathrm{MAP}}(b_3)=a_2.
\]
MAP 的平均正确率为
\[
\frac14+\frac3{20}+\frac18=0.525,
\]
所以
\[
P_{e,\mathrm{MAP}}=1-0.525=0.475.
\]
ML 选中的联合概率为 $1/4,1/10,1/8$，平均正确率为 $0.475$，因此
\[
P_{e,\mathrm{ML}}=1-0.475=0.525.
\]
本题输入不等概，所以 MAP 和 ML 不等价，MAP 的平均错误率更小。
\end{example}

\begin{example}
\textbf{不等概输入下 MAP 与 ML 的差异}：设输入先验概率为
\[
[P_X]=[0.7\quad0.3],
\]
信道转移概率矩阵为
\[
[P_{Y\mid X}]=
\begin{bmatrix}
0.6&0.4\\
0.2&0.8
\end{bmatrix}.
\]
分别求 ML 和 MAP 判决，并比较错误率。

\noindent \textbf{解题流程}：\\[2pt]
步骤一：ML 直接看转移矩阵每一列。\\
步骤二：MAP 先把每行乘以对应先验概率。\\
步骤三：比较两种规则选中的联合概率之和。

\textbf{解}：ML 准则下，第一列 $0.6>0.2$，第二列 $0.8>0.4$，因此
\[
F_{\mathrm{ML}}(b_1)=a_1,
\qquad
F_{\mathrm{ML}}(b_2)=a_2.
\]
联合概率矩阵为
\[
[P_{XY}]=
\begin{bmatrix}
0.42&0.28\\
0.06&0.24
\end{bmatrix}.
\]
MAP 准则逐列比较联合概率，得到
\[
F_{\mathrm{MAP}}(b_1)=a_1,
\qquad
F_{\mathrm{MAP}}(b_2)=a_1.
\]
ML 的平均正确率为 $0.42+0.24=0.66$，所以 $P_{e,\mathrm{ML}}=0.34$；MAP 的平均正确率为 $0.42+0.28=0.70$，所以
\[
P_{e,\mathrm{MAP}}=0.30.
\]
这个例子说明：当某个输入先验概率明显更大时，MAP 可能会选择与 ML 不同的判决，并得到更小的平均错误率。
\end{example}

\subsection{线性分组码与汉明距离}

线性分组码把连续的信息位按块处理。它通过线性关系产生校验位，使所有合法码字形成一个有结构的集合。接收端可以利用“接收序列离哪个合法码字更近”来译码。

\begin{definitionBox}{线性分组码}
把信源输出的信息序列每 $k$ 个信息位分为一段，通过线性编码规则产生 $r$ 个校验位，输出长度为
\[
n=k+r
\]
的码字，称为 $(n,k)$ 线性分组码。二元 $(n,k)$ 线性分组码共有
\[
M=2^k
\]
个码字，码率为 $k/n$。
\end{definitionBox}

如果码字的开头或结尾 $k$ 位直接是信息位，称为系统码；否则称为非系统码。无论是否系统，线性分组码的关键都是码字集合的距离结构。

为了理解“线性”二字，可以知道三个最基本的对象：信息向量 $\mathbf{u}$ 通过生成矩阵 $G$ 产生码字
\[
\mathbf{c}=\mathbf{u}G;
\]
校验矩阵 $H$ 用来描述合法码字集合，合法码字满足
\[
H\mathbf{c}^{\mathrm T}=0;
\]
若接收序列为 $\mathbf{r}=\mathbf{c}+\mathbf{e}$，则伴随式
\[
\mathbf{s}=H\mathbf{r}^{\mathrm T}=H\mathbf{e}^{\mathrm T}
\]
只与错误图样有关。本章不展开生成矩阵、校验矩阵和伴随式译码表的构造；这里只需要把它们看成描述码字集合和错误规律的工具，后续主要使用汉明距离和最小距离判断纠错能力。

\begin{definitionBox}{汉明距离与最小距离}
对两个二元码字
\[
x=[x_1,\dots,x_n],
\qquad
y=[y_1,\dots,y_n],
\]
汉明距离定义为
\[
d(x,y)=\sum_{i=1}^{n}x_i\oplus y_i,
\]
即两个码字对应位置不同的个数。

码字集合的最小距离为
\[
d_{\min}=\min_{i\ne j}d(v_i,v_j).
\]
\end{definitionBox}

汉明距离满足非负性、对称性和三角不等式。对线性分组码而言，任意两个码字的模二和仍是码字，因此最小距离也可以通过非零码字的最小重量计算：
\[
d_{\min}=\min_{i\ne j}d(v_i,v_j)
=\min_{v_k\ne0}w(v_k),
\]
其中 $w(v_k)$ 表示码字中 $1$ 的个数。

\begin{theoremBox}{定理7.1：分组码纠错能力}
最小距离为 $d_H$ 的二元分组码能纠正 $t$ 个错误的充要条件是
\[
d_H\ge 2t+1.
\]
\end{theoremBox}

这个条件的直觉是：若要纠正 $t$ 个错误，每个码字周围半径为 $t$ 的接收球不能互相重叠。因此两个合法码字至少要相距 $2t+1$，译码器才能唯一判断接收序列更接近哪个码字。

\begin{example}
\textbf{最小距离与纠错能力}：设
\[
C_1=\{00000,01010,10101,11111\},
\qquad
C_2=\{00000,11111\}.
\]
分别求两个码的最小距离和纠错能力。

\textbf{解}：对 $C_1$，非零码字最小重量为
\[
w(01010)=2,
\]
所以 $d_{\min}=2$。因为 $2<2\times1+1$，它不能纠正 $1$ 个错误，只能说无法纠错。

对 $C_2$，两个码字的距离为
\[
d(00000,11111)=5,
\]
所以 $d_{\min}=5$。由
\[
5=2\times2+1
\]
可知该码能纠正 $2$ 个错误。
\end{example}

\subsection{序列最大似然译码}

单符号判决只看一个输出符号；真正的信道编码通常按码字序列译码。序列最大似然译码直接比较整个接收序列由各个候选码字产生的概率。

\begin{definitionBox}{序列最大似然译码准则}
若对所有 $i$ 都满足
\[
p(\mathbf{y}\mid \mathbf{x}=\mathbf{c}^*)
\ge
p(\mathbf{y}\mid \mathbf{x}=\mathbf{c}_i),
\]
则译码为码字 $\mathbf{c}^*$ 对应的消息，即
\[
g(\mathbf{y})=f^{-1}(\mathbf{c}^*).
\]
\end{definitionBox}

对无记忆二元对称信道，设交叉错误概率为 $p$。若接收序列 $\mathbf{y}$ 与候选码字 $\mathbf{x}$ 的汉明距离为 $d$，则
\[
p(\mathbf{y}\mid\mathbf{x})=(1-p)^{n-d}p^d.
\]
取对数得
\[
\log p(\mathbf{y}\mid\mathbf{x})
=n\log(1-p)+d_H\log\frac{p}{1-p}.
\]
当 $p\le1/2$ 时，$\log\frac{p}{1-p}\le0$，所以汉明距离越小，似然越大。

\begin{theoremBox}{定理7.2：BSC 上的 ML 与最小距离译码}
对于无记忆二元对称信道，若错误概率 $p\le1/2$，则最大似然译码准则等价于最小汉明距离译码准则。
\end{theoremBox}

这一定理把概率判决转化成距离判决：在 BSC 中，不必逐项比较复杂似然，只要选择与接收序列汉明距离最小的合法码字即可。

\begin{example}
\textbf{BSC 上的序列 ML 译码}：设码字集合为
\[
C=\{00000,01010,10101,11111\},
\]
接收序列为 $\mathbf{y}=00100$，信道为错误概率 $p=0.1$ 的 BSC。求 ML 译码结果。

\textbf{解}：因为 $p=0.1\le1/2$，BSC 上的 ML 译码等价于最小汉明距离译码。分别计算距离：
\[
d(00100,00000)=1,
\quad
d(00100,01010)=4,
\]
\[
d(00100,10101)=2,
\quad
d(00100,11111)=4.
\]
最小距离为 $1$，对应码字 $00000$，因此 ML 译码输出 $00000$ 对应的消息。

若出现两个或多个码字到接收序列的距离同为最小，则 ML 判决出现并列，需要题目额外规定固定的择一规则；否则只能说明这些码字具有相同最大似然。
\end{example}

\subsection{有噪信道编码定理}

前面的 MAP、ML 和分组码分析说明了如何降低错误率；有噪信道编码定理则回答一个更根本的问题：在有噪信道中，可靠传输到底是否可能，以及信息传输速率能达到多高。

以 BSC 为例可以形成一个直觉：长度为 $n$ 的码字经过错误概率为 $p$ 的信道后，典型错误个数大约是 $pn$。接收端需要在每个码字周围留出一团“可能由噪声造成的接收序列”，这团典型噪声范围的规模约为 $2^{nH(p)}$；而全部二进制序列空间大小为 $2^n$。因此能够互不混淆地放入的码字数量大约不超过
\[
\frac{2^n}{2^{nH(p)}}=2^{n(1-H(p))},
\]
对应的可靠传输速率上限就是
\[
C=1-H(p).
\]
这只是帮助理解的典型噪声球直觉，不是定理证明，也不说明有限码长时可以做到零错误。

\begin{theoremBox}{定理7.5：有噪信道编码定理}
设离散无记忆平稳信道的容量为 $C$。若信息传输速率满足
\[
R<C,
\]
则当码长 $n$ 足够大时，总存在一种 $(M,n)$ 码，使译码错误概率 $p_E$ 任意小。

反之，若
\[
R>C,
\]
则对任何编码方式，译码差错率都不能趋于 $0$。
\end{theoremBox}

这个定理也称为香农第二定理。它的意义在于证明了“高可靠性”和“接近容量的高效率”可以同时存在，但它本身是存在性结论，并不直接告诉我们具体应该怎样构造最佳码。Turbo 码、LDPC 码和 Polar 码等都是接近 Shannon 极限的重要编码方案。

\begin{theoremBox}{定理7.6：信源信道编码定理}
设离散无记忆平稳信道的每秒容量为 $C$，离散信源每秒的熵为 $H$。若
\[
H<C,
\]
则总存在一种编码系统，使信源输出能以任意小的错误概率通过信道传输；若
\[
H>C,
\]
则对任何编码系统，译码差错率都大于 $0$。
\end{theoremBox}

信源信道编码定理把信源压缩和信道传输联系起来：信源每秒产生的不确定性不能超过信道每秒可靠承载的信息量。做题时通常就是比较 $H$ 和 $C$，或比较实际传输速率和信道容量。

\begin{example}
\textbf{BSC 容量与传输时间}：有一个二元对称信道，错误率 $p=0.02$，该信道以 $1500$ 个二元符号/秒的速率传送消息。现有一条由 $0,1$ 等概产生、长度为 $14000$ 的二元符号消息序列。问能否在 $10$ 秒内无差错传输？实现无差错传输的最短时间是多少？

\noindent \textbf{解题流程}：\\[2pt]
步骤一：用 BSC 容量公式计算每个信道符号可可靠承载的信息量。\\
步骤二：乘以信道符号速率，得到每秒容量。\\
步骤三：比较 $10$ 秒容量与消息信息量，并求最短时间。

\textbf{解}：二元对称信道容量为
\[
C=1-H(0.02)=0.8596\text{ bit/符号}.
\]
每秒容量为
\[
C_t=0.8596\times1500=1.288\times10^3\text{ bit/s}.
\]
在 $10$ 秒内最多可靠传输
\[
1.288\times10^3\times10<14000
\]
比特，因此不能在 $10$ 秒内无差错传输。

若传输时间为 $T$，要满足
\[
\frac{14000}{T}\le1.288\times10^3,
\]
所以
\[
T\ge\frac{14000}{1.288\times10^3}\approx10.87\text{ s}.
\]
最短时间约为 $10.87$ 秒。
\end{example}

\subsection{题型整理：第七章常见计算}

\textbf{题型快速判断}：看到题目后先判断它要比较概率、距离还是速率。
\begin{itemize}
  \item 给先验概率和转移矩阵 $\to$ 优先考虑 MAP，先求联合概率矩阵，再按列选最大；
  \item 输入等概或先验未知 $\to$ 常用 ML，直接在转移概率矩阵每列选最大；
  \item 给若干码字 $\to$ 计算两两汉明距离，或对线性码找非零码字最小重量；
  \item 问能纠几个错 $\to$ 用 $d_H\ge2t+1$；
  \item BSC 上做序列译码 $\to$ 当 $p\le1/2$ 时选汉明距离最小的码字，若距离并列要说明需要择一规则；
  \item 给 BSC 错误率和传输时间 $\to$ 先算 $C=1-H(p)$，再比较实际速率和容量。
\end{itemize}

\textbf{答题模板}：
\begin{enumerate}
  \item \textbf{MAP/ML 判决题}：写清每个输出对应矩阵的一列；MAP 比较联合概率，ML 比较转移概率；最后用被选中联合概率之和算正确率。
  \item \textbf{分组码距离题}：先列码字集合，再算 $d_{\min}$；若是线性码，可直接找非零码字最小重量。
  \item \textbf{纠错能力题}：由 $d_H\ge2t+1$ 解出最大整数 $t$，不要把检测能力和纠错能力混淆。
  \item \textbf{编码定理题}：把题目给出的实际信息速率写成 $R$ 或 $H$，与信道容量 $C$ 比较；小于容量才有可靠传输的存在性保证，不代表某个有限码长下零错误。
\end{enumerate}

\begin{warningBox}{第七章易错点汇总}
\begin{itemize}
  \item MAP 比较后验概率或联合概率，ML 比较转移概率；两者只有在输入等概等条件下才等价；
  \item $R=\log M/n$，对二进制 $(n,k)$ 线性分组码才可直接写成 $R=k/n$；
  \item 最小距离是码字集合中任意两个不同码字之间的最小距离，不是接收序列到某个码字的距离；
  \item 能纠正 $t$ 个错误的条件是 $d_H\ge2t+1$，不是 $d_H\ge t+1$；
  \item 有噪信道编码定理给出可靠编码的存在性，不等于已经给出了具体编码构造；
  \item $R<C$ 表示随着码长增大错误概率可以任意小，不是说有限长度码一定无差错；
  \item BSC 中 ML 等价于最小汉明距离需要 $p\le1/2$，距离并列时还需要额外的择一规则。
\end{itemize}
\end{warningBox}

\subsection{公式速查表}

{\scriptsize
\begin{tabularx}{\textwidth}{@{}l>{\raggedright\arraybackslash}X@{}}
\toprule
\textbf{类别} & \textbf{公式 / 核心结论} \\
\midrule
信息传输速率 & \textcolor{Red}{$R=\dfrac{\log M}{n}$}，等概消息时每个码符号携带的信息量 \\
\midrule
二进制 $(n,k)$ 码率 & \textcolor{Red}{$R=\dfrac{k}{n}$}，其中 $M=2^k$ \\
\midrule
条件错误率 & $\sum_{a_i\ne a^*}P_{X\mid Y}(x=a_i\mid y=b_j)$ \\
\midrule
平均错误率 & \textcolor{Red}{$P_E=1-\sum_y p(x^*,y)$} \\
\midrule
MAP 准则 & \textcolor{Red}{$p(x=a^*\mid y)\ge p(x=a_i\mid y)$}，等价于逐列比较联合概率 \\
\midrule
ML 准则 & \textcolor{Red}{$p(y\mid x=a^*)\ge p(y\mid x=a_i)$}，输入等概时与 MAP 等价 \\
\midrule
汉明距离 & \textcolor{Red}{$d(x,y)=\sum_{i=1}^{n}x_i\oplus y_i$} \\
\midrule
最小距离 & \textcolor{Red}{$d_{\min}=\min_{i\ne j}d(v_i,v_j)$}；线性码中 $d_{\min}=\min_{v\ne0}w(v)$ \\
\midrule
纠错条件 & \textcolor{Red}{$d_H\ge2t+1$}，可纠正 $t$ 个错误 \\
\midrule
BSC 序列似然 & $p(\mathbf{y}\mid\mathbf{x})=(1-p)^{n-d}p^d$ \\
\midrule
BSC 最佳译码 & \textcolor{Red}{$p\le1/2$ 时，ML 译码等价于最小汉明距离译码} \\
\midrule
BSC 容量直觉 & $2^n/2^{nH(p)}=2^{n(1-H(p))}$，对应 $C=1-H(p)$ \\
\midrule
有噪信道编码定理 & \textcolor{Red}{$R<C$ 时存在码使 $p_E$ 任意小；$R>C$ 时不能可靠传输} \\
\midrule
信源信道编码定理 & \textcolor{Red}{$H<C$ 时可靠传输存在；$H>C$ 时差错率不能趋于 $0$} \\
\bottomrule
\end{tabularx}
}

\subsection{本章要求}
本章学习重点：
\begin{itemize}
  \item 理解信道编码和信道译码通过冗余提高可靠性的思想；
  \item 掌握信息传输速率 $R=\log M/n$ 和二进制 $(n,k)$ 码率 $R=k/n$；
  \item 会写单符号判决规则，并计算条件错误率、平均错误率和平均正确率；
  \item 会根据先验概率和转移概率矩阵使用 MAP 准则；
  \item 理解 ML 准则，并知道输入等概时 ML 与 MAP 等价；
  \item 会完成例7.2 类型的 MAP/ML 判决和错误率计算；
  \item 了解生成矩阵、校验矩阵和伴随式在描述线性分组码时的基本作用；
  \item 掌握线性分组码、系统码、汉明距离、最小距离和码字重量的概念；
  \item 会用 $d_H\ge2t+1$ 判断分组码纠错能力；
  \item 理解 BSC 上序列 ML 译码等价于最小汉明距离译码的条件，并会处理简单序列译码题；
  \item 能解释有噪信道编码定理和信源信道编码定理的含义；
  \item 能用典型噪声球直觉理解 BSC 容量 $C=1-H(p)$ 的来源；
  \item 会用 BSC 容量公式解决传输速率和最短时间问题。
\end{itemize}
